\documentclass[11pt]{article}

\usepackage{graphicx}
\usepackage{framed}
\usepackage{hyperref}
\usepackage{listings}
\usepackage{xcolor}
\usepackage{caption}
\usepackage{subcaption}

\usepackage{booktabs}
\usepackage{amsmath}

\usepackage[backend=biber,style=authoryear]{biblatex}
\addbibresource{references.bib} 

\marginparwidth 0.5in 
\oddsidemargin 0.25in 
\evensidemargin 0.25in 
\marginparsep 0.25in
\topmargin 0.25in 
\textwidth=6in
\textheight=8in

\definecolor{codegreen}{rgb}{0,0.6,0}
\definecolor{codegray}{rgb}{0.5,0.5,0.5}
\definecolor{codepurple}{rgb}{0.58,0,0.82}
\definecolor{backcolor}{rgb}{0.95,0.95,0.95}

% lstlisting setttings
\lstset{
  backgroundcolor=\color{backcolor}, % Set background color
  commentstyle=\color{codegreen}, % Style for comments
  keywordstyle=\color{magenta}, % Style for keywords
  numberstyle=\tiny\color{codegray}, % Style for line numbers
  stringstyle=\color{codepurple}, % Style for strings
  basicstyle=\ttfamily\footnotesize, % Basic font style and size
  breakatwhitespace=false, % Don't break lines at whitespace only
  breaklines=true, % Enable line breaking
  captionpos=b, % Caption position (bottom)
  keepspaces=true, % Keep spaces
  numbers=left, % Show line numbers on the left
  numbersep=5pt, % Separation of numbers from code
  showspaces=false, % Don't show spaces as visible characters
  showstringspaces=false, % Don't show spaces in strings
  showtabs=false, % Don't show tabs as visible characters
  tabsize=2, % Tab size
  language=Java% Specify the language
}

\begin{document}
\hfill\vbox{\hbox{Jude Shin}
		\hbox{CSC 369, Section 01}	
		\hbox{Assignment 1}	
		\hbox{\today}}\par

\bigskip
\centerline{\Large\bf Map/Reduce}\par
\bigskip

\begin{enumerate}
	\item Consider a Map/Reduce solution to a program that reads a file full with 
		integers and computes the number of integers that are divisible by 3. Each 
		line can contain multiple integers. Example input:

		2 324 234 234 234 

		21432 324 3243 2143214

		3214324 24 223 454


		\begin{enumerate}
			\item What is the input/output type of the mapper and reducer?

				\textit{Input key-value type for mapper: (LongWritable, Text). 
				Output key-value type for mapper: (NullWritable, IntWritable). 
				Input key-value type for reducer: (NullWritable, IntWritable). 
				Output key-value type for reducer: (NullWritable, IntWritable).}
				
				\textit{The input for the mapper would be the line number, and the 
				characters that are on the line. The output of the mapper would be 
				null, and the number (count) of integers that are divisible by 3.
				The input and output for the reducer will be the exact same as the 
				output for the mapper. (null, and the count of integers divisible by 
				3).}

			\item Show the map method for the program.
				
\begin{lstlisting}
public class MyMapper 
	extends Mapper<LongWritable, Text, NullWritable, IntWritable> {

	@Override
	public void map(LongWritable key, Text value, Context context) 
		throws IOException, InterruptedException {

		// The string in one line of the input file
		String fullLine = value.toString().trim();

		// Split the numbers by spaces
		String[] tokens = fullLine.split(" ");

		int count = 0;	// The number of integers that we "accepted"
		for (String chunk : tokens) {
			int num = Integer.parseInt(chunk); // The integer version of the string
			if (num%3 == 0) { // Only include those that are divisible by 3
				count++;
			}
		}
		
		// Only write once after summing everything
		context.write(
				NullWritable.get(),			// Nothing...
				new IntWritable(count));	// The number of valid integers in this line
	}
}
\end{lstlisting}

			\item Show the reduce method for the program.
				
\begin{lstlisting}
public class MyReducer 
	extends Reducer<NullWritable, IntWritable, NullWritable, IntWritable> {

	@Override
	public void reduce(NullWritable nw, Iterable<IntWritable> counts, Context context) 
		throws IOException, InterruptedException {
	
		// Add together all of the counts
		int total = 0;
		for (IntWritable c : counts) {
			total+=c.get();
		}

		context.write(NullWritable.get(), new IntWritable(total));
	}
}
\end{lstlisting}

			\item Can you create a combiner for the program? What will be the code 
				for the reduce method of the combiner? Do you need to change the map 
				or reduce method to support the combiner? If so, show the new map and 
				reduce method.
				
				\textit{I wrote my mapper and reducer for compatibility with a 
				combiner. It also just so happens that I can use the reducer as the 
				combiner because we are doing just a simple enumeration over particular 
				integers. Since the reducer has the same inputs and outputs as the 
				output for the mapper, we can reuse the same code.}

		\end{enumerate}

	\item Consider a Map/Reduce solution to a program that reads a file that 
		contains the date and a temperature that is measured for this day on each 
		line. There can be multiple temperatures that are measured each day. The 
		program should calculate the highest temperature for each day. Example 
		input:

		2022/12/10  10

		2022/12/10  8

		2022/12/10  6

		2022/12/12  11

		2000/01/01  3


		\begin{enumerate}
			\item What is the input/output type of the mapper and reducer?
				
				\textit{Input key-value type for mapper: (LongWritable, Text). 
				Output key-value type for mapper: (Text, IntWritable). 
				Input key-value type for reducer: (Text, IntWritable). 
				Output key-value type for reducer: (Text, IntWritable).}

				\textit{The input for the mapper would be the line number, and the 
				characters that are on the line. The output of the mapper would be 
				the date, and the temperature for that day. The input and output for 
				the reducer will be the exact same as the output for the mapper: 
				(date, temperature).}


			\item Show the map method for the program.
				
\begin{lstlisting}
\end{lstlisting}

			\item Show the reduce method for the program.
				
\begin{lstlisting}
\end{lstlisting}

			\item Can you create a combiner for the program? What will be the code 
				for the reduce method of the combiner? Do you need to change the map 
				or reduce method to support the combiner? If so, show the new map and 
				reduce method.
				
				\textit{}

		\end{enumerate}
\end{enumerate}

% Any code/terminal commands can be blocked like so:
% \begin{lstlisting}
% $ yay -Syu #performs a full system upgrade to keep software up to date
% 
% int main(int argc, const char* argv[]) {
% 	fprintf(stdout, "Hello World!\n");
% 	exit(EXIT_SUCCESS);
% }
% \end{lstlisting}
% 
% If you want some inline code blocks, then you can use the lst package with the same settings \lstinline{like so}.
% 
% Quoted stuff is ``weird".

\end{document}
